\section{Discussion and limitations}
\label{sec:discussion}

\paragraph{Where the price comes from.}
The breakdown is unambiguous. Federation in the narrow sense, that is splitting the data
and averaging local models, costs little: under 2 SHD for $K=8$ on average, and it is
explained by client drift when clients hold fewer time points than there are regressors.
Incomplete optimisation is negligible once the order is fixed. Communication delays act
through the number of rounds, and through staleness at $d=20$. Almost the entire price is
\emph{combinatorial}: \psidag's greedy, data-free projection picks a suboptimal
topological order and rarely revises it. With the true order, the same least-squares
estimator is nearly exact. The question ``how much worse is the federated first-order method
than the exact MIP?'' therefore reduces, to first order, to ``how good is the projection's
order?''. Federation is not the issue.

\paragraph{Improving the order.}
Since the ordering term dominates, the most promising improvements of \feddbn{} target the
projection. Options include several restarts selected by the federated loss, which is
embarrassingly parallel and costs one scalar per client per restart; local search over
adjacent transpositions, scored by clients; and projections that account for the scale of
the variables. Any of these must keep the projection free of data, or else pay for it in
communication.

\paragraph{Limitations.}
(i) All clients are homogeneous in distribution, since they hold windows of one stationary series.
Heterogeneous clients (different noise levels or mechanisms) cause client drift
\citep{karimireddy2020scaffold}, and \feddbn{} would then need drift correction.
(ii) We report the SHD at the oracle threshold, as \exdbn{} does. It favours continuous methods,
which do not fix the support themselves. We also store SHD at a fixed threshold of $0.1$ in
the released results.
(iii) \exdbn{} was run with the 120\,s limit of its main experiments, on two threads.
\Cref{tab:budget} shows that its accuracy for $d\ge 20$ is very sensitive to this budget. A
systematic sweep over time limits and threads, and the resulting time to accuracy, is left
for future work. Its worst-case running time is exponential in any case.
(iv) \feddbn{} shares iterates. Iterates of a least-squares problem reveal information about
the clients' Gram matrices, and formal privacy, for example differentially private
aggregation, is future work.
(v) We provide no convergence theory for the federated projection scheme. \citet{ziu2024psi}
analyse the centralised stochastic-approximation iteration. Extending this analysis to local
steps and bounded delays is open.

\paragraph{Conclusion.}
We proposed \feddbn, a federated extension of \psidag{} to dynamic Bayesian networks, and
measured its price of federation against the exact MIP \exdbn{}, which by
\cref{prop:gram} can itself be federated with one round of secure aggregation. The price
is real for small models: about 27\% of the true edges for $d\le 15$. Almost all of it is
due to the greedy order of the projection, not to splitting the data, incomplete
optimisation or communication delays. For larger models the comparison turns around in
practice, even though the MIP remains better in principle. With 120\,s per instance the MIP
is no better than \feddbn{} at $d=20$--$25$. It needs several times more compute to recover
the true graph there, and from $d=100$ on it cannot solve even its root relaxation. \feddbn{}
keeps an SHD of about 0.4 of the true edges up to $d=1000$, at a cost that grows polynomially
and is known in advance. Better data-free projections are the natural next step: they
target the one term that dominates the price.
